%-------------------------------------------------------------------
\chapter{Conclusions}
\label{c:conc}
%-------------------------------------------------------------------

\epigraph{\enquote{Computer Science is no more about computers than astronomy is about telescopes.}}{\emph{E. W. Dijkstra}}

This chapter summarizes the conclusions of the work and outlines limitations and possible extensions.

\section{Summary and limitations}
\label{sec:conc:summary}

The thesis showed that the spatial neighborhood size \(NB\) in Mixture NCA has a measurable impact on long-horizon tissue-like dynamics: \(NB=1\) is inadequate, while \(NB \ge 2\) yields a clear improvement; beyond that, the ``best'' \(NB\) depends on the metric and the model (Mixture vs.\ Stochastic Mixture).
The main threats to validity are summarized in \cref{sec:impl:threats}: dependence on the synthetic simulator and dataset size, the limited range of \(NB\) explored, and the absence of validation-based model selection in the curriculum experiment.
These caveats do not invalidate the core finding---that neighborhood size matters and that a controlled protocol is needed to observe it---but they bound the generality of the conclusions.

\section{Promising directions for future work}
\label{sec:conc:future}

Several directions could strengthen or extend this work:
\begin{itemize}
  \item \textbf{Longer training on long windows:} train for more epochs on 500-step windows (or add further curriculum phases) to reduce the gap between generated and ground-truth trajectories at long horizons.
  \item \textbf{Validation loss and early stopping:} compute a validation loss using free rollouts and use it for early stopping, so that optimization is aligned with the evaluation regime.
  \item \textbf{Final-state and trajectory-level losses:} add losses that directly reward how close the final state of a rollout is to the ground truth (e.g., distribution of cell types), or that penalize drift over the full trajectory.
  \item \textbf{Scheduled teacher forcing:} gradually reduce the fraction of steps where the true state is fed during training, so that the model learns to correct its own errors and exposure bias is mitigated.
  \item \textbf{Evaluation with multiple rollouts:} for the stochastic model, run several rollouts per initial condition and report averaged or best-of-\(k\) metrics to better capture typical behavior.
\end{itemize}
We leave these ideas for future work; the current experiments already establish that the neighborhood-size effect is real and carries over to the long-horizon curriculum setting.
